\subsection{Producto de espacios vectoriales}

\begin{frame}
    \frametitle{Un ejemplo antes de la definición}

    Supongamos que queremos manejar a la vez un vector de $\mathbb{R}^2$
    y un polinomio de $\mathbb{R}_1[x]$, es decir, pares
    \[
        \bigl((a,b),\,c+dx\bigr),
        \qquad
        (a,b)\in\mathbb{R}^2,
        \quad
        c+dx\in\mathbb{R}_1[x].
    \]

    Es natural sumar y multiplicar por escalares
    \alert{componente a componente}:
    \[
        \bigl((1,2),\,3-x\bigr)+\bigl((0,-1),\,1+4x\bigr)
        =
        \bigl((1,1),\,4+3x\bigr),
    \]
    \[
        2\,\bigl((1,2),\,3-x\bigr)
        =
        \bigl((2,4),\,6-2x\bigr).
    \]

    \medskip

    Cada componente se opera en su propio espacio vectorial,
    con independencia de la otra.
\end{frame}


\begin{frame}
    \frametitle{Producto de dos espacios vectoriales}

    \begin{definition}
        Sean $V$ y $W$ espacios vectoriales sobre el mismo cuerpo $k$.
        El \alert{producto (cartesiano)} de $V$ y $W$ es el conjunto
        \[
            V\times W=\{(v,w):v\in V,\ w\in W\}
        \]
        con las operaciones
        \[
            (v,w)+(v',w')=(v+v',\,w+w'),
            \qquad
            \lambda(v,w)=(\lambda v,\,\lambda w),
        \]
        para $v,v'\in V$, $w,w'\in W$ y $\lambda\in k$.
    \end{definition}

    \medskip

    En la primera componente se usan las operaciones de $V$,
    y en la segunda, las de $W$.
\end{frame}


\begin{frame}[allowframebreaks]
    \frametitle{El producto es un espacio vectorial}

    \begin{proposition}
        Con las operaciones anteriores, $V\times W$ es un espacio vectorial
        sobre $k$. Su vector cero es $(0,0)$ y el opuesto de $(v,w)$ es
        $(-v,-w)$.
    \end{proposition}

    \begin{proof}
        Hay que comprobar los axiomas. La idea es la misma en todos:
        cada axioma en $V\times W$ se reduce al mismo axioma en $V$
        (primera componente) y en $W$ (segunda componente).

        Por ejemplo, la distributividad respecto de la suma de vectores:
        \begin{align*}
            \lambda\bigl((v,w)+(v',w')\bigr)
            &=\lambda(v+v',\,w+w')
            =\bigl(\lambda(v+v'),\,\lambda(w+w')\bigr)\\
            &=(\lambda v+\lambda v',\,\lambda w+\lambda w')
            =(\lambda v,\lambda w)+(\lambda v',\lambda w')\\
            &=\lambda(v,w)+\lambda(v',w'),
        \end{align*}
        donde en la segunda línea se han usado los axiomas en $V$ y en $W$.

        Para el vector cero,
        \[
            (v,w)+(0,0)=(v+0,\,w+0)=(v,w),
        \]
        y para el opuesto,
        \[
            (v,w)+(-v,-w)=\bigl(v+(-v),\,w+(-w)\bigr)=(0,0).
        \]

        Los axiomas restantes se comprueban del mismo modo, componente
        a componente.
    \end{proof}
\end{frame}


\begin{frame}
    \frametitle{Combinaciones lineales en el producto}

    Como las operaciones se hacen componente a componente, lo mismo
    ocurre con las combinaciones lineales.

    \begin{proposition}
        Para $(v_1,w_1),\ldots,(v_r,w_r)\in V\times W$ y
        $\lambda_1,\ldots,\lambda_r\in k$,
        \[
            \lambda_1(v_1,w_1)+\cdots+\lambda_r(v_r,w_r)
            =
            \bigl(\lambda_1v_1+\cdots+\lambda_rv_r,\;
                  \lambda_1w_1+\cdots+\lambda_rw_r\bigr).
        \]
    \end{proposition}

    \begin{proof}
        Por la definición del producto por escalares,
        $\lambda_i(v_i,w_i)=(\lambda_iv_i,\lambda_iw_i)$ para cada $i$.
        Sumando estos $r$ pares y aplicando repetidamente la definición
        de la suma en $V\times W$, se obtiene el par cuyas componentes son
        las sumas de las componentes.
    \end{proof}

    \medskip

    En particular, una combinación lineal en $V\times W$ es nula
    si y solo si \alert{ambas} componentes lo son.
\end{frame}


\begin{frame}
    \frametitle{Ejemplo: $k^m\times k^n$ y $k^{m+n}$}

    \begin{example}
        Eliminando los paréntesis interiores, a cada par de
        $k^m\times k^n$ le corresponde un vector de $k^{m+n}$:
        \[
            \bigl((a_1,\ldots,a_m),\,(b_1,\ldots,b_n)\bigr)
            \;\longleftrightarrow\;
            (a_1,\ldots,a_m,b_1,\ldots,b_n).
        \]
        Esta correspondencia es biyectiva y respeta las operaciones:
        sumar o multiplicar por escalares antes o después de eliminar
        los paréntesis da el mismo resultado, porque en ambos casos
        se opera coordenada a coordenada.
    \end{example}

    \medskip

    Por eso, en la práctica, identificamos
    \[
        k^m\times k^n=k^{m+n}.
    \]
    En el Tema~4 daremos un sentido preciso a esta identificación.
\end{frame}


\begin{frame}
    \frametitle{Producto de varios espacios vectoriales}

    \begin{definition}
        Sean $V_1,\ldots,V_r$ espacios vectoriales sobre $k$.
        Su \alert{producto} es
        \[
            V_1\times\cdots\times V_r
            =
            \{(u_1,\ldots,u_r):u_i\in V_i\},
        \]
        con la suma y el producto por escalares componente a componente:
        \[
            (u_1,\ldots,u_r)+(u'_1,\ldots,u'_r)=(u_1+u'_1,\ldots,u_r+u'_r),
        \]
        \[
            \lambda(u_1,\ldots,u_r)=(\lambda u_1,\ldots,\lambda u_r).
        \]
    \end{definition}

    Como en el caso de dos factores, es un espacio vectorial.
    El caso más conocido es $k^n=k\times\cdots\times k$, con $n$ factores.
\end{frame}


\begin{frame}
    \frametitle{Una base del producto: un ejemplo}

    En $\mathbb{R}^2\times\mathbb{R}_1[x]$, todo par se descompone como
    \[
        \bigl((a,b),\,c+dx\bigr)
        =
        \bigl((a,b),\,0\bigr)+\bigl((0,0),\,c+dx\bigr),
    \]
    y por tanto
    \[
        \bigl((a,b),\,c+dx\bigr)
        =
        a\bigl((1,0),0\bigr)+b\bigl((0,1),0\bigr)
        +c\bigl((0,0),1\bigr)+d\bigl((0,0),x\bigr).
    \]

    \medskip

    Los cuatro vectores que aparecen se obtienen
    \begin{itemize}
        \item de la base canónica de $\mathbb{R}^2$, poniendo $0$
              en la segunda componente;
        \item de la base $\{1,x\}$ de $\mathbb{R}_1[x]$, poniendo $(0,0)$
              en la primera componente.
    \end{itemize}

    \medskip

    \begin{center}
        \alert{¿Forman siempre una base del producto los vectores así obtenidos?}
    \end{center}
\end{frame}


\begin{frame}
    \frametitle{Bases del producto}

    \begin{theorem}
        Sean $\{v_1,\ldots,v_m\}$ una base de $V$ y
        $\{w_1,\ldots,w_n\}$ una base de $W$. Entonces
        \[
            B=\{(v_1,0),\ldots,(v_m,0),\,(0,w_1),\ldots,(0,w_n)\}
        \]
        es una base de $V\times W$. En particular,
        \[
            \dim(V\times W)=\dim V+\dim W.
        \]
    \end{theorem}

    \medskip

    Hay que probar dos cosas:
    \begin{enumerate}
        \item $B$ genera $V\times W$;
        \item $B$ es linealmente independiente.
    \end{enumerate}
    En ambas usaremos que las combinaciones lineales en $V\times W$
    se calculan componente a componente.
\end{frame}


\begin{frame}[allowframebreaks, t]
    \frametitle{Bases del producto: demostración}

    \begin{proof}
        \alert{$B$ genera $V\times W$.}
        Sea $(v,w)\in V\times W$. Como $\{v_1,\ldots,v_m\}$ genera $V$ y
        $\{w_1,\ldots,w_n\}$ genera $W$, existen escalares tales que
        \[
            v=\lambda_1v_1+\cdots+\lambda_mv_m,
            \qquad
            w=\mu_1w_1+\cdots+\mu_nw_n.
        \]
        Calculando componente a componente,
        \begin{align*}
            &\lambda_1(v_1,0)+\cdots+\lambda_m(v_m,0)
            +\mu_1(0,w_1)+\cdots+\mu_n(0,w_n)\\
            &\qquad=
            \bigl(\lambda_1v_1+\cdots+\lambda_mv_m,\;
                  \mu_1w_1+\cdots+\mu_nw_n\bigr)
            =(v,w),
        \end{align*}
        ya que en la primera componente los sumandos $\mu_j\cdot0$ son
        nulos, y en la segunda lo son los $\lambda_i\cdot0$.

        \framebreak

        \alert{$B$ es linealmente independiente.}
        Supongamos que
        \[
            \lambda_1(v_1,0)+\cdots+\lambda_m(v_m,0)
            +\mu_1(0,w_1)+\cdots+\mu_n(0,w_n)=(0,0).
        \]
        El mismo cálculo de antes muestra que el miembro izquierdo es
        \[
            \bigl(\lambda_1v_1+\cdots+\lambda_mv_m,\;
                  \mu_1w_1+\cdots+\mu_nw_n\bigr).
        \]
        Igualando componentes,
        \[
            \lambda_1v_1+\cdots+\lambda_mv_m=0
            \quad\text{en } V,
            \qquad
            \mu_1w_1+\cdots+\mu_nw_n=0
            \quad\text{en } W.
        \]
        Como $v_1,\ldots,v_m$ son linealmente independientes,
        $\lambda_1=\cdots=\lambda_m=0$; y como $w_1,\ldots,w_n$ lo son,
        $\mu_1=\cdots=\mu_n=0$.

        Por tanto $B$ es una base de $V\times W$. Tiene $m+n$ vectores
        (son distintos dos a dos, por serlo los $v_i$ y los $w_j$, y por
        no ser nulos), luego
        \[
            \dim(V\times W)=m+n=\dim V+\dim W.\qedhere
        \]
    \end{proof}
\end{frame}


\begin{frame}
    \frametitle{Coordenadas en la base del producto}

    \begin{proposition}
        Sean $B_V=\{v_1,\ldots,v_m\}$ una base de $V$,
        $B_W=\{w_1,\ldots,w_n\}$ una base de $W$ y
        \[
            B=\{(v_1,0),\ldots,(v_m,0),\,(0,w_1),\ldots,(0,w_n)\}.
        \]
        Para todo $(u,u')\in V\times W$,
        \[
            (u,u')_B=
            \begin{pmatrix}
                u_{B_V}\\
                u'_{B_W}
            \end{pmatrix}
            \in k^{m+n}.
        \]
    \end{proposition}


    Es decir, se \alert{apilan} las coordenadas de cada componente
    en su propia base.
\end{frame}


\begin{frame}
    \frametitle{Coordenadas en la base del producto}

    \begin{proof}
        Si $u_{B_V}=(\lambda_1,\ldots,\lambda_m)$ y
        $u'_{B_W}=(\mu_1,\ldots,\mu_n)$, entonces
        \[
            u=\lambda_1v_1+\cdots+\lambda_mv_m,
            \qquad
            u'=\mu_1w_1+\cdots+\mu_nw_n.
        \]
        Calculando componente a componente,
        \begin{align*}
            (u,u')
            &=
            \bigl(\lambda_1v_1+\cdots+\lambda_mv_m,\;
                  \mu_1w_1+\cdots+\mu_nw_n\bigr)\\
            &=\lambda_1(v_1,0)+\cdots+\lambda_m(v_m,0)
            +\mu_1(0,w_1)+\cdots+\mu_n(0,w_n),
        \end{align*}
        luego
        \[
            (u,u')_B=(\lambda_1,\ldots,\lambda_m,\mu_1,\ldots,\mu_n).
            \qedhere
        \]
    \end{proof}
\end{frame}


\begin{frame}
    \frametitle{Dimensión de un producto de varios espacios}

    \begin{corollary}
        Si $V_1,\ldots,V_r$ tienen dimensión finita, entonces
        \[
            \dim(V_1\times\cdots\times V_r)
            =
            \dim V_1+\cdots+\dim V_r.
        \]
    \end{corollary}

    \begin{proof}
        Se elige una base de cada $V_i$ y se coloca cada uno de sus vectores
        en la componente $i$-ésima, con $0$ en las demás. La misma
        demostración que para dos factores, ahora con $r$ componentes,
        prueba que todos estos vectores forman una base del producto.
        Su número es $\dim V_1+\cdots+\dim V_r$.
    \end{proof}

    \medskip

    Por ejemplo, $\dim k^n=\dim k+\cdots+\dim k=n$, y la base obtenida
    es precisamente la base canónica $\{e_1,\ldots,e_n\}$.
\end{frame}


\begin{frame}
    \frametitle{Ejemplos de dimensión de productos}

    \begin{example}
        \begin{itemize}
            \item $\dim\bigl(\mathbb{R}^2\times\mathbb{R}_1[x]\bigr)=2+2=4$.
            \item $\dim\bigl(M_{2\times2}(k)\times k_3[x]\bigr)=4+4=8$.
            \item $\dim\bigl(k^m\times k^n\bigr)=m+n$, de acuerdo con la
                  identificación $k^m\times k^n=k^{m+n}$.
            \item $\dim\bigl(M_{m\times n}(k)\times M_{m\times n}(k)\bigr)=2mn$.
        \end{itemize}
    \end{example}

    En $\mathbb{R}^2\times\mathbb{R}_1[x]$, con la base
    \[
        B=\bigl\{((1,0),0),\ ((0,1),0),\ ((0,0),1),\ ((0,0),x)\bigr\},
    \]
    el vector $\bigl((3,-1),\,2-5x\bigr)$ tiene coordenadas
    \[
        \bigl((3,-1),\,2-5x\bigr)_B
        =
        \begin{pmatrix}3\\-1\\2\\-5\end{pmatrix}.
    \]
\end{frame}



\begin{frame}
    \frametitle{No toda base del producto es de esta forma}

    \begin{example}
        En $\mathbb{R}\times\mathbb{R}=\mathbb{R}^2$, los vectores
        \[
            (1,1),
            \qquad
            (1,-1)
        \]
        forman una base, pues la matriz
        \[
            \begin{pmatrix}
                1&1\\
                1&-1
            \end{pmatrix}
        \]
        tiene rango $2$.

        Sin embargo, ninguno de ellos es de la forma $(v,0)$ ni $(0,w)$.
    \end{example}

    \medskip

    El teorema proporciona \alert{una} base de $V\times W$ a partir
    de bases de $V$ y de $W$, pero $V\times W$ tiene muchas otras bases
    que no se obtienen así.
\end{frame}
